% !PRINTED 389
\weberpage{389}

But if $\sigma$ is different from 1, then one can choose a pair of figures $\alpha, \beta$, both among the $1,2,3,4$ or both among the $5,6,7,8$, in such a way that the transposition
\[
\sigma\,(\alpha,\beta)\,\sigma^{-1} = (\alpha', \beta')
\]
is different from $(\alpha, \beta)$. For if one first takes for $\alpha$ a figure changed by $\sigma^{-1}$, then $\alpha'$ is different from $\alpha$ (Bd.\,I, \S.\,154, 6.); if one then chooses for $\beta$ a figure which does not pass into $\alpha$ under $\sigma^{-1}$, and which occurs together with $\alpha$ in the first or in the second half of the eight figures — such a figure certainly always exists, because only one figure passes into $\alpha$ under $\sigma^{-1}$ — then $(\alpha', \beta')$ is different from $(\alpha, \beta)$. But it then follows that in the group $Q$ the element $(\alpha,\beta)\,\sigma\,\Pi_{1,2,3,4}\,(\alpha,\beta)$, and consequently also [by \S.\,102, (8), (9)]
\[
(\alpha,\beta)\,\sigma\,\Pi_{1,2,3,4}\,(\alpha,\beta)\,\Pi_{1,2,3,4}\,\sigma^{-1} = (\alpha,\beta)\,\sigma\,(\alpha,\beta)\,\sigma^{-1} = (\alpha,\beta)(\alpha',\beta')
\]
occurs, and this is a permutation from $S$ different from the identity. With this, then, it is proved:

3.\ \gesperrt{The group $P$ of the double-tangent equation is simple.}

As a special application we may point out that the group $P$ can effect among the 28 roots of the double-tangent equation only permutations of the first kind, and that \gesperrt{consequently the discriminant of this equation is a square}.

For if we denote for the moment by $G$ the group of all permutations of the 28 roots, and by $A$ the alternating group contained in it as a normal divisor, then the greatest common divisor of $A$ and $P$ is a normal divisor of $P$ and must therefore be identical with $P$; that is, $P$ is contained in $A$.

If one seeks out the group of the permutations which leave one of the roots of the double-tangent equation, say $[1\,2]$, unchanged, one finds a group which is still transitive on the remaining 27 roots\footnote{By a geometrical theorem of \gesperrt{Geiser}, this group is isomorphic with the group of the equation of degree 27, on which the solution of the problem of the 27 straight lines on a surface of the third order depends. (Mathem. Annalen, Bd.\,I.).}.
